\documentclass{article}

\usepackage{amsmath,amssymb}
\usepackage{xurl}
\usepackage[hidelinks]{hyperref}
\setlength{\emergencystretch}{2em}

% Shared commands for notes in docs/paper-gaps/.

\DeclareUnicodeCharacter{2080}{\ifmmode{}_{0}\else\textsubscript{0}\fi}
\DeclareUnicodeCharacter{2081}{\ifmmode{}_{1}\else\textsubscript{1}\fi}
\DeclareUnicodeCharacter{2082}{\ifmmode{}_{2}\else\textsubscript{2}\fi}
\DeclareUnicodeCharacter{2099}{\ifmmode{}_{n}\else\textsubscript{n}\fi}

\newcommand{\Lean}{\textsc{Lean}}
\ExplSyntaxOn
\NewDocumentCommand{\leanid}{m}{
  \ifmmode
    \textsf{\detokenize{#1}}
  \else
    \group_begin:
    \urlstyle{sf}
    \tl_set:Nn \l_tmpa_tl {#1}
    \tl_replace_all:Nnn \l_tmpa_tl {\_} {_}
    \exp_args:NV \path \l_tmpa_tl
    \group_end:
  \fi
}
\ExplSyntaxOff
\providecommand{\tr}{\operatorname{tr}}
\newcommand{\tnleanIssueBase}{https://github.com/LionSR/TNLean/issues/}
\newcommand{\tnleanPrBase}{https://github.com/LionSR/TNLean/pull/}
\newcommand{\tnleanDocBase}{https://sirui-lu.com/TNLean/docs/}
\newcommand{\ghissue}[1]{\href{\tnleanIssueBase#1}{GitHub issue \##1}}
\newcommand{\ghpr}[1]{\href{\tnleanPrBase#1}{PR \##1}}
\newcommand{\doclink}[1]{\href{\tnleanDocBase#1.html}{\path{#1}}}

% Dirac notation and the identity operator, as in blueprint/src/macros/common.tex.
% \providecommand keeps note-local definitions authoritative.
\usepackage{dsfont}
\providecommand{\ket}[1]{|#1\rangle}
\providecommand{\bra}[1]{\langle #1|}
\providecommand{\braket}[2]{\langle #1 | #2 \rangle}
\providecommand{\ketbra}[2]{|#1\rangle\!\langle #2|}
\providecommand{\Id}{\mathds{1}}
\providecommand{\one}{\mathds{1}}
\providecommand{\C}{\mathbb{C}}
\providecommand{\MN}[1]{M_{#1}(\C)}
\providecommand{\MD}{M_D(\C)}

% External referencing for paper sources.  The build helper writes sanitized
% auxiliary files under build/paper-aux/ and excludes duplicated source labels.
\usepackage{xr}

% Import a generated paper auxiliary file with a caller-supplied label prefix.
% The candidate paths support builds from docs/paper-gaps/, the repository root,
% and build/paper-aux/ itself.
\newcommand{\importpaperaux}[2]{%
  \IfFileExists{../../build/paper-aux/#2.aux}{%
    \externaldocument[#1]{../../build/paper-aux/#2}%
  }{%
    \IfFileExists{build/paper-aux/#2.aux}{%
      \externaldocument[#1]{build/paper-aux/#2}%
    }{%
      \IfFileExists{#2.aux}{%
        \externaldocument[#1]{#2}%
      }{}%
    }%
  }%
}

% General external reference macros
\newcommand{\paperref}[2]{\ref{#1-#2}}
\newcommand{\papereqref}[2]{(\ref{#1-#2})}

% CPSV16 source (1606.00608 / MPDO-22-12-17-2.tex) external document and shortcuts
\importpaperaux{cpsv16-}{1606.00608/MPDO-22-12-17-2-tnlean}
\newcommand{\cpsvref}[1]{\paperref{cpsv16}{#1}}
\newcommand{\cpsveqref}[1]{\papereqref{cpsv16}{#1}}

% Verdict marker: \gapnote{<kind>}{<status>}. Kind is the policy
% classification (clarification, local-correction, scope-restriction,
% unfaithful, false-source, open-gap); status is open, wip (elimination
% underway), resolved, or historical. Machine-read by the site index;
% prints a verdict line.
\newcommand{\gapnote}[2]{\par\noindent\textbf{Verdict:} #1 (#2).\par}


\title{The One-Step Subspace in the Quantum Wielandt Inequality}
\author{}
\date{2026-08-22}

\begin{document}
\maketitle
\gapnote{scope-restriction}{open}

\section{The assertion}

The quantum Wielandt inequality appears as Theorem~1 of
Ref.~\cite{Sanz2010Quantum}.\footnote{For traceability, this note corresponds
    to \href{https://github.com/LionSR/TNLean/issues/1049}{issue \#1049}. The
    local source passages are
    \path{Papers/0909.5347/main.tex:360--365},
    \path{Papers/0909.5347/main.tex:572--590}, and
    \path{Papers/0909.5347/main.tex:645--684}.}
Let $\mathcal E_A$ be a quantum channel on $\MD$ with Kraus operators
$A_1,\ldots,A_a$. For each positive integer $n$, define the exact-length word
span
\begin{align}
    S_n(A)
        &=
        \operatorname{span}
        \left\{A_{i_1}\cdots A_{i_n}
            : 1\leq i_1,\ldots,i_n\leq a\right\}.
    \label{eq:spwc10_wielandt_one_step_subspace_exact_word_span}
\end{align}
Thus $S_n(A)$ is generated by words of length exactly $n$. When a positive
length with full word span exists, define
\begin{align}
    i(A)
        &=
        \min\left\{n\geq 1:S_n(A)=\MD\right\}.
    \label{eq:spwc10_wielandt_one_step_subspace_i_index}
\end{align}
The formal convention discussed below assigns $i(A)=0$ when this set is empty.

For a nonzero vector $\varphi\in\C^D$, let
$H_n(A,\varphi)=S_n(A)\varphi$. With the positive word-length convention of
Ref.~\cite{Sanz2010Quantum}, the primitivity index is
\begin{align}
    q(\mathcal E_A)
        &=
        \min\left\{q\geq 1:
            H_q(A,\varphi)=\C^D
            \text{ for every }\varphi\neq 0\right\}.
    \label{eq:spwc10_wielandt_one_step_subspace_q_index}
\end{align}
This index is likewise assigned the value zero when the defining set is empty.

The source writes $d$ for the number of linearly independent Kraus operators
in a minimal spanning Kraus list. Equivalently,
\begin{align}
    d
        &=
        \dim S_1(A).
    \label{eq:spwc10_wielandt_one_step_subspace_source_kraus_count}
\end{align}
Under this convention, Theorem~1 of Ref.~\cite{Sanz2010Quantum} states that a
primitive channel satisfies
\begin{align}
    i(A)
        &\leq
        (D^2-d+1)D^2.
    \label{eq:spwc10_wielandt_one_step_subspace_general_source_bound}
\end{align}
The same theorem gives two sharper bounds under additional hypotheses on
$S_1(A)$. If $S_1(A)$ contains an invertible matrix, then
\begin{align}
    i(A)
        &\leq
        D^2-d+1.
    \label{eq:spwc10_wielandt_one_step_subspace_invertible_source_bound}
\end{align}
If $S_1(A)$ contains a singular matrix with a nonzero eigenvalue, then
\begin{align}
    i(A)
        &\leq
        D^2.
    \label{eq:spwc10_wielandt_one_step_subspace_singular_source_bound}
\end{align}
Both hypotheses concern arbitrary elements of $S_1(A)$. The source does not
require the matrix with the stated property to occur as a Kraus operator in a
chosen representation.

\section{What has been formalized}

The formal development uses the same exact-length spaces and defines the Kraus
rank by
\begin{align}
    \operatorname{kr}(A)
        &=
        \dim S_1(A).
    \label{eq:spwc10_wielandt_one_step_subspace_formal_kraus_rank}
\end{align}
By (\ref{eq:spwc10_wielandt_one_step_subspace_source_kraus_count}) and
(\ref{eq:spwc10_wielandt_one_step_subspace_formal_kraus_rank}),
$\operatorname{kr}(A)$ is the invariant form of the source parameter $d$.
Replacing $d$ by $\operatorname{kr}(A)$ therefore does not change the
bound.\footnote{The formal definitions are
    \leanid{Kraus.wordSpan} in
    \path{TNLean/Wielandt/SpanGrowth/CumulativeSpan.lean:57--64} and
    \leanid{MPSTensor.krausRank} in
    \path{TNLean/Wielandt/Primitivity/Definitions.lean:84--90}. The local
    source identifies $d$ as the number of linearly independent Kraus
    operators in \path{Papers/0909.5347/main.tex:360--365}; the proof of
    Lemma~1 uses $\dim T_1(A)=d$ in
    \path{Papers/0909.5347/main.tex:579--583}.}

The formal general inequality is
\begin{align}
    i(A)
        &\leq
        (D^2-\operatorname{kr}(A)+1)D^2.
    \label{eq:spwc10_wielandt_one_step_subspace_general_formal_bound}
\end{align}
Its proof passes through eventual full Kraus rank and a blocking argument, but
its conclusion asserts $S_n(A)=\MD$ at one fixed length. It is not merely a
statement about an increasing sum of word spans.\footnote{The general theorem
    is \leanid{MPSTensor.iIndex_le_general_of_isPrimitivePaper} in
    \path{TNLean/Wielandt/Inequality/Bounds.lean:389--468}. The formal
    equivalence, under normalization, between the primitive condition of
    Ref.~\cite{Sanz2010Quantum} and eventual full Kraus rank is in
    \path{TNLean/Wielandt/Primitivity/Equivalence.lean:169--191}.}

The improved formal estimates now use the source hypotheses on the one-step
subspace. In the invertible case, the formal theorem assumes an invertible
matrix $X\in S_1(A)$ and proves
\begin{align}
    S_{D^2-\operatorname{kr}(A)+1}(A)
        &=
        \MD.
    \label{eq:spwc10_wielandt_one_step_subspace_invertible_exact_span}
\end{align}
In the singular case, it assumes a singular matrix $X\in S_1(A)$ with a
nonzero eigenvalue and proves
\begin{align}
    S_{D^2}(A)
        &=
        \MD.
    \label{eq:spwc10_wielandt_one_step_subspace_singular_exact_span}
\end{align}
These are precisely the corresponding hypotheses in Theorem~1 of
Ref.~\cite{Sanz2010Quantum}. The older statements for a chosen Kraus operator
remain as corollaries obtained by setting $X=A_{i_0}$.\footnote{The relevant
    declarations are
    \leanid{MPSTensor.wordSpan_eq_top_of_isPrimitivePaper_of_mem_wordSpan_one_of_isUnit}
    and
    \leanid{MPSTensor.iIndex_le_of_mem_wordSpan_one_of_isUnit} for the
    invertible case, and
    \leanid{MPSTensor.wordSpan_eq_top_of_isPrimitivePaper_of_mem_wordSpan_one_of_noninvertible_eigenvector}
    and
    \leanid{MPSTensor.iIndex_le_sq_of_mem_wordSpan_one_of_noninvertible_eigenvector}
    for the singular case, all in
    \path{TNLean/Wielandt/Inequality/Bounds.lean}. The shorter
    chosen-generator statements in the same file are corollaries.}

\section{Resolved one-step-subspace issue}

Let $X\in S_1(A)$ be invertible. Because $X$ is a linear combination of the
Kraus operators, multiplication on either side by $X$ maps an exact-length word
span into the next exact-length word span. In particular,
\begin{align}
    S_n(A)X
        &\subseteq
        S_{n+1}(A).
    \label{eq:spwc10_wielandt_one_step_subspace_multiplication_inclusion}
\end{align}
Multiplication by $X$ is injective on $\MD$, so
(\ref{eq:spwc10_wielandt_one_step_subspace_multiplication_inclusion}) supplies
the dimension-growth step behind the sharper bound: the dimensions of the
spaces $S_n(A)$ cannot stagnate for too many successive lengths.

The formal proof adjoins the chosen matrix $X\in S_1(A)$ as a redundant
generator. It then proves that the exact word spans and the Kraus rank remain
unchanged. The relevant declarations are
\leanid{Kraus.wordSpan_oneStepAugment_eq} and
\leanid{MPSTensor.krausRank_oneStepAugment}, both in
\path{TNLean/Wielandt/SpanGrowth/InvertibleWordSpan.lean}. This construction
reduces the source hypothesis to the single-generator argument without
changing the channel invariant in
(\ref{eq:spwc10_wielandt_one_step_subspace_general_formal_bound}).

The singular case uses the same construction. If $X\in S_1(A)$ is singular
but has a nonzero eigenvalue, adjoining $X$ preserves the exact word spans and
the Kraus rank. The resulting theorem therefore applies to any such element of
$S_1(A)$, not only to a distinguished Kraus matrix.

\section{Primitive-to-normal theorem with a chosen positive fixed point}

The theorem
\leanid{MPSTensor.isNormal_of_isPrimitiveMPS_with_posDef}\footnote{See
    \path{TNLean/Wielandt/Primitivity/StronglyIrreducibleToFullRank.lean}.}
derives normality from a chosen fixed point $\rho$, the local complementary
transfer-map gap predicate \leanid{MPSTensor.IsPrimitiveMPS}, and the additional
hypothesis $\rho>0$.

This theorem formalizes one route for Proposition~3(c)$\Rightarrow$(b) of
Ref.~\cite{Sanz2010Quantum}. The argument first interprets the assumptions as
strong irreducibility with a positive-definite fixed point and then applies the
eventual-full-Kraus-rank theorem. Positive definiteness enters through the
quadratic form
\begin{align}
    Q_\rho(B)
        &=
        \operatorname{Re}\tr(B^\dagger\rho B).
    \label{eq:spwc10_wielandt_one_step_subspace_trace_pairing}
\end{align}
When $\rho>0$, the form in
(\ref{eq:spwc10_wielandt_one_step_subspace_trace_pairing}) is nondegenerate,
and compactness of the unit sphere gives a uniform positive lower bound. If
only $\rho\geq 0$ is known, the form can vanish on the kernel of $\rho$.

The hypothesis $\rho>0$ is therefore a genuine assumption of this auxiliary
route, not an accidental artefact. The auxiliary theorem should not be cited
as the source theorem itself. The formal source-facing results remain the
primitive-channel statements in
\path{TNLean/Wielandt/Inequality/Bounds.lean} and the
primitive/eventual-full-rank equivalence in
\path{TNLean/Wielandt/Primitivity/Equivalence.lean}. Section~2.3 of
Ref.~\cite{Cirac2016MPDO_arXiv} should be connected to those
primitive-channel statements rather than to the auxiliary theorem requiring a
chosen positive-definite fixed point.

\section{What is not a discrepancy}

The identity
(\ref{eq:spwc10_wielandt_one_step_subspace_formal_kraus_rank}) is neither a
source error nor a formal weakening. It makes explicit the linearly independent
Kraus count used in Ref.~\cite{Sanz2010Quantum}.

Nor does the use of cumulative spaces by itself create a theorem-level
discrepancy. Define
\begin{align}
    T_n(A)
        &=
        S_0(A)+S_1(A)+\cdots+S_n(A).
    \label{eq:spwc10_wielandt_one_step_subspace_cumulative_span}
\end{align}
Some intermediate formal results prove
\begin{align}
    T_{D^2}(A)
        &=
        \MD
    \label{eq:spwc10_wielandt_one_step_subspace_cumulative_full_span}
\end{align}
under normality. The cumulative conclusion
(\ref{eq:spwc10_wielandt_one_step_subspace_cumulative_full_span}) is weaker
than the full quantum Wielandt theorem. The final theorem in
\path{TNLean/Wielandt/Inequality/Bounds.lean}, however, proves the bound for
$i(A)$ defined by the exact-length spaces in
(\ref{eq:spwc10_wielandt_one_step_subspace_i_index}).\footnote{The
    cumulative-span theorem is
    \leanid{Kraus.cumulativeSpan_eq_top} in
    \path{TNLean/Wielandt/SpanGrowth/NonzeroTraceProduct.lean:164--197}; the
    exact-length theorem is the general result identified above. The blueprint
    distinguishes these statements in
    \path{blueprint/src/chapter/ch08_wielandt.tex:461--490} and states the
    general bound in
    \path{blueprint/src/chapter/ch08_wielandt.tex:492--548}.}

The intermediate use of normality is also legitimate. In these formal
statements, normality means eventual full Kraus rank. Under the normalization
used in Ref.~\cite{Sanz2010Quantum}, the development proves that the source
primitive condition is equivalent to eventual full Kraus rank.

\section{Scope of the index comparison}

For $H_n(A,\varphi)=S_n(A)\varphi$, Proposition~1 of
Ref.~\cite{Sanz2010Quantum} states
\begin{align}
    q(\mathcal E_A)
        &\leq
        i(A)
    \label{eq:spwc10_wielandt_one_step_subspace_source_index_comparison}
\end{align}
for every quantum channel.

The available formal comparison assumes trace-preserving normalization and
uniform vector-spreading primitivity. It proves
(\ref{eq:spwc10_wielandt_one_step_subspace_source_index_comparison}) under
those assumptions, so it specializes Proposition~1 rather than reproducing
its full hypothesis set. A companion comparison instead assumes that $i(A)$
is finite.\footnote{The two declarations are
    \leanid{MPSTensor.qIndex_le_iIndex_of_isPrimitivePaper} and
    \leanid{MPSTensor.qIndex_le_iIndex}. The definitions are in
    \path{TNLean/Wielandt/Primitivity/Definitions.lean}, and the restricted
    theorem is in \path{TNLean/Wielandt/Inequality/Bounds.lean}.}

The omitted cases are degenerate but still require an explicit formal
argument. If the channel is not primitive, then the admissible set in
(\ref{eq:spwc10_wielandt_one_step_subspace_q_index}) is empty and
$q(\mathcal E_A)=0$. If the channel lacks eventual full Kraus rank, then the
admissible set in
(\ref{eq:spwc10_wielandt_one_step_subspace_i_index}) is empty and $i(A)=0$.
A source-faithful theorem can therefore prove these two zero-index
characterizations, split on primitivity, and invoke the existing
primitive-channel inequality in the nontrivial case.

\section{Scope of the strong-irreducibility equivalence}

Proposition~3(c) of Ref.~\cite{Sanz2010Quantum} requires that $1$ be the only
peripheral eigenvalue and that its eigenspace be one-dimensional, generated by
a positive-definite fixed point. The one-dimensionality condition is
essential. The identity channel has peripheral spectrum $\{1\}$ and admits
the positive-definite fixed point $I$, but its fixed space is not
one-dimensional.

The present formal predicate strengthens the source condition by also
requiring irreducibility of the transfer map. The proved equivalence is
therefore a strengthened version of Proposition~3, not a literal
formalization of its third clause.\footnote{The strengthened predicate is
    \leanid{MPSTensor.IsStronglyIrreduciblePaper}; the equivalences are
    \leanid{MPSTensor.primitivePaper_iff_stronglyIrreducible} and
    \leanid{MPSTensor.hasEventuallyFullKrausRank_iff_stronglyIrreducible}. See
    \path{TNLean/Wielandt/Primitivity/Definitions.lean} and
    \path{docs/glossary.md}.}

Removing this restriction requires a literal source predicate consisting of a
positive-definite fixed point, a one-dimensional fixed eigenspace, and no
other peripheral eigenvalues. One can then prove that a normalized completely
positive map with these properties has no nontrivial invariant projection:
such a projection would produce an additional positive fixed direction,
contradicting simplicity of the fixed space. This implication supplies the
irreducibility conjunct of the existing predicate, after which the established
equivalences yield the literal Proposition~3. Alternatively, the three
directions of the equivalence can be proved directly for the source predicate.

\section{Conclusion}

The general quantum Wielandt bound is formalized with the exact-length index
and with the source's linearly independent Kraus count written invariantly.
The two improved estimates now use the source hypotheses: the distinguished
matrix may be any element of $S_1(A)$ with the required spectral property.
Statements for a chosen Kraus operator remain useful corollaries, but they are
not the source theorem.

Two scope restrictions remain. The index comparison assumes primitivity or,
in its companion form, finiteness of $i(A)$. The strong-irreducibility
predicate also contains an irreducibility conjunct absent from the literal
third clause of Proposition~3 in Ref.~\cite{Sanz2010Quantum}. The preceding
sections identify the zero-index and source-predicate arguments needed to
remove these restrictions.

Cumulative-span declarations and bundled analysis chains remain support
lemmas. The source-facing quantum Wielandt conclusions concern the
exact-length spaces $S_n(A)$ and the index $i(A)$.

\bibliographystyle{alpha}
\bibliography{references}

\end{document}
